% Adapted from dictionaries/snapshot.tex using a dictionary comprehension.
\begin{blocksection}
\question Complete \lstinline{snapshot}. It takes a function from integers
to integers and a list of integer inputs. Return a dictionary whose keys are
those inputs and whose values are the corresponding outputs of the function.
Use a dictionary comprehension.

\begin{lstlisting}
from collections.abc import Callable

def snapshot(
    f: Callable[[int], int], snap_inputs: list[int]
) -> dict[int, int]:
    """
    >>> snapshot(lambda x: x**2, [1, 2, 3])
    {1: 1, 2: 4, 3: 9}
    >>> snapshot(lambda x: x + 1, [])
    {}
    """
    return {
        ____________: ____________
        for _______________________________
    }
\end{lstlisting}
\begin{solution}
\begin{lstlisting}
def snapshot(
    f: Callable[[int], int], snap_inputs: list[int]
) -> dict[int, int]:
    return {x: f(x) for x in snap_inputs}
\end{lstlisting}
\end{solution}
\end{blocksection}

\begin{questionmeta}
\textbf{Teaching Tips}\par\nopagebreak[4]
Suggested Time: 7 min; Difficulty: Easy--Medium
\nopagebreak[4]
\begin{itemize}
    \item \textbf{Read the function type:} In \lstinline{Callable[[int], int]},
    the bracketed list gives the parameter types and the final type describes
    the return value. Here \lstinline{f} accepts one integer and returns one.
    \item \textbf{Connect to HOFs:} The argument is a function, not the result
    of calling it. Evaluate \lstinline{f(x)} separately for each input.
    \item \textbf{Choose key and value:} Each input \lstinline{x} is a key;
    its output \lstinline{f(x)} is the value. The comprehension constructs
    the complete result, with no separate dictionary-update step.
    \item \textbf{Explain the name:} The dictionary records a function's
    behavior on a finite selection of inputs: a snapshot of its behavior.
\end{itemize}
\end{questionmeta}
